\section{The \system{} architecture}
\label{sec:method}

\subsection{Relational temporal detector}

The detector organizes the 86 HAI variables as nodes and retains six relation channels instead of collapsing every edge into one adjacency matrix. The channels encode official fluid relations, official measurement relations, device-channel pairing, operational context, dependencies estimated only from normal training windows, and constrained cross-subsystem hardware-in-the-loop dependencies. This provenance separation matters: a learned correlation is useful context, but it is not mislabeled as a verified physical connection.

Let $A_r\in\{0,1\}^{N\times N}$ be relation $r\in\{1,\ldots,R\}$ and let $\widetilde A_r$ denote its normalized form. A scalar measurement and a learned node embedding are projected to $d=32$ dimensions. The relational update used before and after temporal encoding can be represented as
\begin{align}
M_t &= \sum_{r=1}^{R}\widetilde A_r H_t W_r, \\
g_t &= \sigma_g(H_tW_g+b_g), \\
\bar H_t &= \operatorname{LN}\!\left(H_t+g_t\odot\rho(M_t)\right),
\label{eq:relgraph}
\end{align}
where $W_r$ is relation-specific, $g_t$ is a learned gate, $\rho$ is a nonlinearity, and $\operatorname{LN}$ is layer normalization. The checkpoint contains six relation-specific linear maps in both the pre-temporal and post-temporal graph blocks.

Temporal dependence is modeled with one four-head Transformer encoder layer \citep{vaswani2017attention}. For head $h$,
\begin{equation}
\operatorname{Attn}_h(Q_h,K_h,V_h)=
\operatorname{softmax}\!\left(\frac{Q_hK_h^\top}{\sqrt{d_h}}\right)V_h,
\label{eq:attention}
\end{equation}
and the concatenated heads are followed by a feed-forward block, residual connections, and normalization. Learned time embeddings cover the 64-row window. Attention pooling produces a window representation $z$, while node representations support a root-cause head. The frozen checkpoint exposes an attack head, root head, and state decoder:
\begin{align}
p_a &= \sigma(f_{\mathrm{attack}}(z)), \\
\boldsymbol{\pi} &= \operatorname{softmax}\!\left(f_{\mathrm{root}}(H_{1:T})/\tau_r\right), \\
\widehat X &= f_{\mathrm{dec}}(H_{1:T}).
\label{eq:heads}
\end{align}
The configuration uses five epochs, batch size eight, learning rate $2\times10^{-4}$, weight decay $10^{-4}$, dropout 0.15, and normal-only detection mode. Its auxiliary reconstruction and root coefficients are 0.05 and 0.25. The configured explicit physics-invariant weight is zero and the invariant list is empty; accordingly, we describe the checkpoint as \emph{process-graph grounded}, not physics-informed. This distinction is retained in the limitations.

An alarm is issued at $p_a\geq0.5$. The top-ranked root probabilities provide causal candidates, but they are interpreted as model support rather than identified physical causality. The model is intentionally compact so the full gate can run on CPU-scale infrastructure.

\subsection{Closed evidence bundle}

For an alarmed window, the evidence assembler records the alarm probability, ranked root support, node values, two alternative causes, uncertainty, process context, provenance status, and a small permitted-action set. Every evidence item receives a stable identifier, value, reliability field, score, and critical flag. The bundle also records whether the predicted node appears in the configured topology, how many invariants were expected and evaluated, and the result of any configured counterfactual support check.

The bundle is \emph{closed}: generation may select or restate supplied content but may not create new telemetry, causes, times, actions, or certainty. A SHA-256 window digest records post-ingestion bytes. It does not authenticate the acquisition source, so acquisition-boundary provenance is marked unavailable rather than overstated.

\subsection{Canonical operator renderer}

The hardened design makes integrity-critical presentation a deterministic function $C$ of the evidence bundle:
\begin{equation}
C(\mathcal B)=\bigl(s^*,I^*,A^*,q^*,r^*\bigr).
\label{eq:canonical}
\end{equation}
The renderer sorts evidence by support, includes every critical item, and fills the surface to at least three items when available. It names at most two alternatives. When uncertainty is at least 0.25, the operator question asks for context that distinguishes the predicted cause from the leading alternative; otherwise it asks whether authorized maintenance explains the event. The recommended action is the first item in the bundle's permitted diagnostic list. Alarm wording is selected directly from the thresholded probability.

The generator supports a deterministic template and an optional Hugging Face mode. In the latter, an LLM can propose a structured object, but the summary, evidence order, alternatives, question, and recommendation are overwritten with $C(\mathcal B)$ before verification. The stated cause and uncertainty remain independently checked. The external LLM used for red teaming is separate from this generator and has no authority over the bundle.

The post-freeze utility study adds a finite registered-surface extension. Let $\mathcal T_j$ be a preregistered rendering template and $\theta_j(\mathcal B)$ its bundle-derived slots. The admissible set is
\begin{equation}
\mathcal R(\mathcal B)=\{C(\mathcal B)\}\cup
\bigl\{\mathcal T_j(\theta_j(\mathcal B)):j\in\mathcal J_{\mathrm{safe}}\bigr\}.
\label{eq:registry}
\end{equation}
The registry is finite, versioned, and evaluated after the primary held-out freeze. It permits only renderings whose operator-critical fields and slots were registered in advance; it is not an open-ended semantic paraphrase classifier.

\subsection{Independent decision-semantic gate}

The verifier receives both $\mathcal B$ and $x$. It checks identifier membership, critical-evidence coverage, cause agreement, uncertainty range and tolerance, absence of absolute-certainty language, alternative completeness, permitted action, topology support, invariants, counterfactual support, and provenance. Table~\ref{tab:verifier} lists the contract.

\begin{table}[t]
\centering
\caption{Decision-semantic verification contract. Hard failures produce \textsc{abstain}; integrity violations that do not meet a hard-failure rule produce \textsc{warning}.}
\label{tab:verifier}
\small
\begin{tabularx}{\linewidth}{@{}>{\raggedright\arraybackslash}p{0.28\linewidth}>{\raggedright\arraybackslash}Xl@{}}
\toprule
Check & Acceptance condition & Failure class \\
\midrule
Evidence identifiers & Every identifier exists; all critical evidence is present & hard / warning \\
Cause & Stated cause belongs to the bundle and equals the predicted cause & hard / warning \\
Uncertainty & Value is in $[0,1]$ and within tolerance 0.02 of the bundle & hard / warning \\
Alternatives & Required top-two alternatives are present; no unsupported cause appears & hard / warning \\
Action & Recommendation exactly equals one permitted diagnostic action & hard \\
Surface commitment & Summary, evidence order, alternatives, question, and action equal the canonical rendering & hard \\
Unsafe framing & No downplaying, delay, false-alarm, or routine-continuation language & warning \\
Grounding & Topology supports the cause and required invariants are evaluated & warning \\
Counterfactual & Support meets the configured minimum & warning \\
Provenance & Required acquisition-boundary status is verified & hard \\
\bottomrule
\end{tabularx}
\end{table}


The primary hardened configuration adds three decision-semantic checks. First, every field in $\Psi(x)$ must equal $C(\mathcal B)$. Second, a deny-list catches explicit downplaying or delay language as defense in depth. Third, critical evidence may not be moved below weaker noncritical evidence. In the registered extension, the equality test becomes $\Psi(x)\in\mathcal R(\mathcal B)$ while every fact, order, action, topology, and provenance check remains active. If $c_j\in[0,1]$ denotes the $j$th soft check, the diagnostic score is
\begin{equation}
S(\mathcal B,x)=\frac{1}{J}\sum_{j=1}^{J}c_j.
\label{eq:verifier_score}
\end{equation}
Hard failures override the score and return \abstain. Integrity violations without a hard failure return \warning. Otherwise, scores at least 0.80 return \verified, scores from 0.55 to 0.80 return \warning, and smaller scores return \abstain. The gate therefore exposes reasons rather than a silent binary decision.

The legacy comparison uses the same evidence checks but disables decision-semantic enforcement. In particular, it accepts recommendations when a permitted action appears as a substring and does not require equality with the canonical surface. This comparison isolates the effect of binding what the operator sees, rather than replacing the whole detection pipeline.

\subsection{Contract-level acceptance invariant}

The canonical design yields a simple bounded guarantee.

\begin{proposition}[Unregistered-surface rejection]
Fix a bundle $\mathcal B$ and a finite registered set $\mathcal R(\mathcal B)$. If the hardened verifier returns \verified{} only when $\Psi(x)\in\mathcal R(\mathcal B)$, then every candidate $x'$ with $\Psi(x')\notin\mathcal R(\mathcal B)$ fails to receive \verified{} status. The canonical verifier is the special case $\mathcal R(\mathcal B)=\{C(\mathcal B)\}$.
\label{prop:surface}
\end{proposition}

\noindent The proof follows by substitution and contradiction and is given in Appendix~\ref{app:proof}. This is a contract property, not a universal natural-language safety theorem. It depends on a complete definition of the integrity-critical surface and on uncompromised bundle, renderer, and verifier implementations.

\subsection{Freeze and traceability}

The development red team was evaluated against the legacy verifier. Its observed failures informed canonicalization. The verifier, prompts, task files, model choice, decoding settings, and relevant artifact hashes were then refrozen before held-out responses were generated. At evaluation time, the script recomputes the verifier and task-file hashes and refuses to continue on drift. The later multi-model and registered-surface analyses are explicitly marked post-freeze development studies and do not alter the primary held-out implementation or claim. This sequence prevents tuning on held-out outcomes and makes every reported table and figure traceable to machine-readable inputs.

\begin{table}[t]
\centering
\caption{Publication-freeze identifiers and disclosure boundaries. Full machine-readable values are included in the Overleaf package.}
\label{tab:reproducibility}
\small
\begin{tabularx}{\linewidth}{@{}p{0.31\linewidth}X@{}}
\toprule
Item & Frozen identifier or status \\
\midrule
Freeze schema & \texttt{coact-cps-publication-freeze-v1} \\
Freeze ID & \texttt{60b263ab81ea\ldots e1034bf22} \\
Verifier SHA-256 & \texttt{38c4255cdaca\ldots 80151a6dc041} \\
Detector checkpoint SHA-256 & \texttt{8bc18c69992e\ldots 5d2e468f7f8d} \\
Adjacency SHA-256 & \texttt{4d32889436e3\ldots 02bba96912455} \\
Held-out tasks SHA-256 & \texttt{9c9759e360ce\ldots b66839cdb416e} \\
Restricted prepared test bytes & Excluded from the evidence archive; checksum recorded, not redistributed \\
Ground-truth labels in red-team tasks & Not exported \\
Human semantic review & Finalized on the development sample; operator study still required \\
\bottomrule
\end{tabularx}
\end{table}

